\section{Key Technical Ideas}
\label{sec:ideas}

The proof is organized around obstacles that recur across operations, rather
than around the order of source files or tactics.

\subsection{Relating idealized and bounded-exponent arithmetic}

Earlier analyses express significand rounding in Flocq's FLX format, whose
exponent range is unbounded~\cite{Rideau2022}.  Binary64 instead has finite
upper and lower exponent bounds and gradual underflow.  A bounded-format proof
therefore follows one of two routes.  When an intermediate lies outside the
underflow region, a bridge lemma identifies its bounded FLT rounding with the
corresponding FLX rounding, allowing the existing analysis to transfer at that
node.  When a designated intermediate can be subnormal or zero, the proof
reasons directly about FLT rounding and carries the resulting absolute error.

This distinction is local to rounded nodes.  It does not justify one global
rule that every operation must reject subnormal values, nor does the phrase
``whole range'' mean that no other hypotheses are needed.  Finiteness,
overflow headroom, input representation, and nonzero-result hypotheses remain
operation-specific.

\subsection{Mixed bounds near underflow}

The standard relative model becomes uninformative when an exact result
approaches zero, because a fixed absolute rounding quantum cannot always be
absorbed into a uniform relative factor.  The running example therefore states
both a relative-scale term and an absolute underflow term in
\eqref{eq:dw-times-fp-mixed}.  The decisive proof step is to track which residual
may be lost and to bound its contribution to the represented result, while
separately proving that the remaining computation stays finite.

Other operations make different choices: in some cases the same coefficient
survives gradual underflow; in others the whole-range result has a different
coefficient or a mixed form.  The library interface preserves these distinctions
instead of forcing every operation into one generic error statement.

\subsection{Conditions as an interface between proofs}

An operation condition packages the facts expected by the numerical theorem in
a form that the VST development can name in a function specification.  A bridge
lemma proves that the condition implies the lower-level hypotheses of the
accuracy and representation lemmas.  Independently, a check proof establishes
that a successful Boolean result implies the condition.  The condition is thus
the stable boundary between numerical reasoning and program verification.

This separation also permits multiple interfaces to the same kernel.  One
guard may enforce conditions for a stronger exactness or representation
result, while another may accept a larger domain and expose only an accuracy
property.  The theorem used by the wrapper, rather than the implementation name
alone, determines the caller-visible guarantee.

\subsection{Accuracy and closure}

The real-number error bound and the double-word relation solve different
client problems.  Accuracy controls the represented value.  Closure provides
the representation fact expected at a subsequent call.  In algorithms ending
with Fast2Sum, the closure proof typically depends on showing that the final
high/low reconstruction preserves the intended intermediate sum under the
applicable range and ordering facts.  The paper states these properties
separately so that a contract does not silently treat one as a consequence of
the other.
